%% %%%%%%%%%%%%%%%%%%%%%%%%%%%%%%%%%%%%%%%%%%%%%%%%%%%%%%%%%%%%%%%
%% AGFA-Book-Kapitel-1.tex -- Kapitel I
%% %%%%%%%%%%%%%%%%%%%%%%%%%%%%%%%%%%%%%%%%%%%%%%%%%%%%%%%%%%%%%%%
\chapter{Grundlagen}
\label{chap:grundlagen}

Ein kurzer Überblick über das Kapitel steht vor der ersten Section.

\section{Metrische Räume}
\label{sec:metrisch}

\begin{definition}
\label{def:metrik}
Eine Abbildung $d\colon X\times X\to\mathbb{R}$ heißt \emph{Metrik}, wenn
für alle $x,y,z\in X$ gilt
\begin{equation}
\label{eq:dreieck}
  d(x,z)\le d(x,y)+d(y,z).
\end{equation}
\end{definition}

\begin{example}
Der Betrag liefert die Metrik $d(x,y)=\abs{x-y}$ auf $\mathbb{R}$.
\end{example}

\section{Vollständigkeit}
\label{sec:vollstaendig}

\begin{lemma}
\label{lem:cauchy}
Jede konvergente Folge ist eine Cauchy-Folge.
\end{lemma}

\begin{proof}
Folgt aus der Dreiecksungleichung \eqref{eq:dreieck}.
\end{proof}

\subsection{Kontraktionen}
\label{subsec:kontraktion}

\begin{theorem}[Banach]
\label{thm:banach}
Ist $(X,d)$ vollständig und $T\colon X\to X$ eine Kontraktion, so hat $T$
genau einen Fixpunkt.
\end{theorem}

\begin{corollary}
\label{cor:eindeutig}
Die Fixpunktiteration konvergiert für jeden Startwert.
\end{corollary}

Die Abschätzung
\begin{equation}
\label{eq:apriori}
  d(x_n,x^*)\le\frac{q^n}{1-q}\,d(x_1,x_0)
\end{equation}
zeigt die Geschwindigkeit.
